\chapter{Il sistema endocrino}

% ---------------------------------------------------------------------------
\section{Generalità}

\textbf{Sistema endocrino}: comprende tutte le cellule e i tessuti endocrini
presenti nel corpo umano.

\begin{itemize}
  \item le \textbf{cellule endocrine} sono cellule ghiandolari che riversano
    gli ormoni da esse prodotti nel liquido interstiziale, nell'apparato
    linfoide e nel sangue;
  \item le \textbf{cellule esocrine}, al contrario, rilasciano le proprie
    secrezioni su una superficie epiteliale.
\end{itemize}

% ---------------------------------------------------------------------------
\section{Classificazione degli ormoni}

\rubrica{Classificazione strutturale}
\begin{enumerate}
  \item \textbf{derivati degli aminoacidi}: molecole di dimensioni ridotte,
    strutturalmente simili agli aminoacidi; es.\ derivati dalla tirosina
    (ormoni tiroidei e catecolamine, adrenalina e noradrenalina), derivati
    del triptofano (melatonina);
  \item \textbf{ormoni peptidici}: formati da catene di aminoacidi,
    rappresentano il gruppo più numeroso; es.\ tutti gli ormoni ipofisari;
  \item \textbf{ormoni steroidei}: derivati del colesterolo; sono gli ormoni
    dell'apparato riproduttivo e della ghiandola surrenale;
  \item \textbf{eicosanoidi}: piccole molecole con un anello a 5 atomi di
    carbonio a un'estremità; coordinano le attività cellulari e influenzano
    processi enzimatici (es.\ la coagulazione del sangue) che si verificano
    nei liquidi extracellulari.
\end{enumerate}

\rubrica{Classificazione biochimica}
\begin{itemize}
  \item \textbf{proteici}: ormone della crescita (GH), prolattina (PRL),
    ormone adrenocorticotropo (ACTH), paratormone (PTH), calcitonina,
    insulina, glucagone, somatostatina, polipeptide pancreatico, ormoni
    ipotalamici (CRH, TRH, GHRH, GnRH), inibenti VIP, vasopressina (ormone
    antidiuretico, ADH), ossitocina;
  \item \textbf{glicoproteici}: ormone tireotropo (TSH), ormone
    follicolo-stimolante (FSH), ormone luteinizzante (LH), gonadotropina
    corionica (hCG);
  \item \textbf{steroidei}: cortisolo, aldosterone, $17\beta$-estradiolo,
    progesterone, testosterone;
  \item \textbf{derivati da amine}: tetraiodotironina (T4), triiodotironina
    (T3), adrenalina, noradrenalina.
\end{itemize}

% ---------------------------------------------------------------------------
\section{Sintesi, trasporto e recettori}

\begin{itemize}
  \item gli ormoni sono normalmente prodotti nel reticolo endoplasmatico
    delle rispettive cellule sotto forma di \textbf{pro-ormoni} e accumulati
    nell'apparato di Golgi, spesso racchiusi in vescicole; possono essere
    attivati all'interno o all'esterno delle cellule;
  \item sono in genere trasportati nel sangue legati a \textbf{specifiche
    proteine trasportatrici} (albumine o globuline), che ne prolungano la
    vita in circolo;
  \item i \textbf{recettori} possono trovarsi sulle membrane cellulari o sui
    nuclei, quando la struttura dell'ormone permette il suo passaggio
    attraverso la membrana;
  \item la \textbf{concentrazione ematica} di ciascun ormone dipende
    dall'equilibrio fra velocità di sintesi (o di liberazione) e velocità di
    eliminazione, ed è costantemente regolata in maniera molto precisa.
\end{itemize}

% ---------------------------------------------------------------------------
\section{L'asse ipotalamo-ipofisi}

Il principale sistema endocrino è l'\textbf{asse ipotalamo-ipofisi}, che
regola la funzione di numerose altre ghiandole endocrine e collega il
sistema endocrino con quello nervoso.

\begin{itemize}
  \item l'ipotalamo è collegato con l'ipofisi attraverso un peduncolo che
    contiene vasi e nervi: i primi formano il \textbf{circolo portale
    ipotalamo-ipofisario}, mentre i secondi realizzano la neurosecrezione
    nell'ipofisi posteriore;
  \item tutti gli ormoni prodotti dall'ipofisi, e in parte anche i fattori
    rilascianti prodotti dall'ipotalamo, agiscono principalmente sulle loro
    cellule bersaglio, ma, essendo ormoni che entrano in circolo, hanno anche
    effetti diretti su altre cellule (non tutti noti);
  \item l'attività dell'asse ipotalamo-ipofisi è prevalentemente controllata
    a \textbf{feedback negativo} dagli stessi ormoni che esso controlla.
\end{itemize}

% ---------------------------------------------------------------------------
\section{I riflessi endocrini}

L'attività endocrina è controllata da riflessi endocrini, che rispondono a:
\begin{enumerate}
  \item \textbf{stimoli umorali}: variazione nella composizione del fluido
    extracellulare;
  \item \textbf{stimoli ormonali}: arrivo o rimozione di uno specifico
    ormone;
  \item \textbf{stimoli nervosi}: arrivo di neurotrasmettitori a livello
    delle giunzioni neuroendocrine.
\end{enumerate}
Sono regolati da forme di feedback, negativo e positivo.

% ---------------------------------------------------------------------------
\section{Panoramica delle ghiandole endocrine}

\begin{itemize}
  \item \textbf{ipotalamo}: produzione di ADH e ossitocina, secrezione di
    ormoni di regolazione;
  \item \textbf{ipofisi}: adenoipofisi (ACTH, TSH, GH, PRL, FSH, LH, MSH),
    neuroipofisi (rilascio di ADH e ossitocina);
  \item \textbf{epifisi}: melatonina;
  \item \textbf{ghiandole paratiroidi} (sulla superficie posteriore della
    tiroide): paratormone;
  \item \textbf{tiroide}: tiroxina (T4), triiodotironina (T3), calcitonina;
  \item \textbf{timo} (prima dell'involuzione nell'adulto): timosina;
  \item \textbf{cuore}: peptide natriuretico atriale, peptide natriuretico
    encefalico;
  \item \textbf{ghiandola surrenale}: midollare (adrenalina, noradrenalina),
    corticale (cortisolo, corticosterone, aldosterone, androgeni);
  \item \textbf{rene}: eritropoietina (EPO), calcitriolo;
  \item \textbf{tessuto adiposo}: leptina, resistina;
  \item \textbf{tratto digerente}: numerosi ormoni;
  \item \textbf{isole pancreatiche}: insulina, glucagone;
  \item \textbf{gonadi}: testicolo (androgeni/testosterone, inibina), ovaio
    (estrogeni/estradiolo, progestinici, inibina).
\end{itemize}

% ---------------------------------------------------------------------------
\section{L'ipofisi}

Localizzata nella \textbf{sella turcica} dello sfenoide, è connessa
all'ipotalamo dall'\textbf{infundibolo}. Si distinguono
\textbf{adenoipofisi} (lobo anteriore, con pars distalis, pars intermedia,
pars tuberalis) e \textbf{neuroipofisi} (lobo posteriore, pars nervosa).

\rubrica{Controllo ipotalamico}
L'ipotalamo esercita tre tipi di controllo sugli organi endocrini:
\begin{enumerate}
  \item secrezione degli ormoni regolatori che controllano l'attività
    dell'adenoipofisi;
  \item produzione diretta di ADH e ossitocina, rilasciati dalla
    neuroipofisi;
  \item innervazione simpatica diretta della midollare del surrene, con
    secrezione di adrenalina e noradrenalina.
\end{enumerate}

\rubrica{Il sistema portale ipofisario}
L'adenoipofisi è connessa all'ipotalamo dal \textbf{circolo portale
ipotalamo-ipofisario}. I sistemi portali forniscono un efficiente mezzo di
comunicazione chimica, assicurando che tutto il sangue che entra nel circolo
portale raggiunga le cellule bersaglio prima di ritornare alla circolazione
generale: la comunicazione è rigidamente a senso unico, poiché tutte le
sostanze chimiche secrete devono compiere un giro completo del sistema
circolatorio prima di raggiungere i capillari all'origine del circolo
portale.

\rubrica{Ormoni della neuroipofisi (lobo posteriore, pars nervosa)}
\begin{itemize}
  \item \textbf{vasopressina} (ADH, ormone antidiuretico): agisce sul
    tubulo contorto distale del rene, con riassorbimento di acqua e aumento
    di volume e pressione del sangue;
  \item \textbf{ossitocina}: agisce su utero e ghiandola mammaria (femmina)
    con contrazione del miometrio ed eiezione di latte, e su dotto
    deferente e prostata (maschio) con contrazione ed eiezione delle
    secrezioni.
\end{itemize}

I nuclei ipotalamici \textbf{sopraottico} e \textbf{paraventricolare}
mandano i loro assoni nella parte posteriore dell'ipofisi. Nelle
terminazioni nervose sono presenti vescicole contenenti i rispettivi
ormoni: quando i neuroni sono stimolati, i potenziali d'azione provocano la
fusione delle vescicole con le membrane cellulari e l'esocitosi delle
molecole di ormone, che entrano nel torrente circolatorio.

% ---------------------------------------------------------------------------
\section{Patologie legate al GH}

\rubrica{Durante l'accrescimento}
\begin{itemize}
  \item \textbf{gigantismo}: da ipersecrezione di GH;
  \item \textbf{nanismo ipofisario}: da iposecrezione di GH, con riduzione
    dell'attività della cartilagine epifisaria e formazione di ossa
    insolitamente corte.
\end{itemize}

\rubrica{Nell'adulto}
\textbf{Acromegalia}: crescita abnorme di ossa piatte (fronte, zigomi),
ingrossamento delle dita, diabete; si sviluppa quando i livelli sierici di
GH aumentano dopo la chiusura delle epifisi, per cui lo scheletro non si
allunga ma le ossa diventano più spesse.

% ---------------------------------------------------------------------------
\section{L'ormone della crescita (GH)}

L'unico ormone ipofisario che agisce direttamente, non attraverso la
regolazione di altri organi endocrini, è l'\textbf{ormone della crescita}
(\textit{Growth Hormone}, GH).

\begin{itemize}
  \item \textbf{ruolo principale}: regolatore della crescita dell'organismo
    e del metabolismo;
  \item \textbf{meccanismo d'azione}: stimola la sintesi proteica, la
    gluconeogenesi (da precursori lipidici) e la lipolisi; provoca
    ritenzione di sodio e di azoto;
  \item azione specifica sulle cartilagini di accrescimento;
  \item \textbf{effetto indiretto}: stimola la produzione di \textbf{IGF}
    (\textit{insuline-like growth factor}) nel fegato e in altri tessuti;
    l'IGF-1 coadiuva le azioni di GH, ma ha effetti opposti sul metabolismo,
    riducendo la glicemia e favorendo la deposizione di trigliceridi.
\end{itemize}

% ---------------------------------------------------------------------------
\section{La prolattina}

L'adenoipofisi produce una serie di ormoni che regolano l'attività
secretoria di altri organi a secrezione interna (tiroide, surrenale,
genitali), tra cui il GH e la \textbf{prolattina} (PRL), il cui effetto
principale consiste nello stimolare la produzione del latte da parte delle
ghiandole mammarie, principalmente dopo il parto e durante la lattazione.

\begin{itemize}
  \item la sua produzione è spontanea, ma normalmente inibita da un fattore
    ipotalamico detto \textbf{PIF} (\textit{prolactin inhibiting factor}),
    oggi identificato con la \textbf{dopamina};
  \item è anche fortemente influenzata dal sistema nervoso centrale, in
    particolare dal lobo limbico, sede delle influenze emotive;
  \item la suzione del capezzolo e lo svuotamento degli alveoli mammari ne
    stimolano la produzione, mentre il rigonfiamento della mammella per
    mancata lattazione la inibisce;
  \item sono noti gli effetti della lattazione (e della sua sospensione)
    sull'umore, probabilmente legati al ruolo della dopamina (la cui
    riduzione provoca depressione).
\end{itemize}

\rubrica{Effetti sulla fertilità}
\begin{itemize}
  \item la prolattina inibisce la produzione di GnRH, il fattore
    ipotalamico che stimola la secrezione di gonadotropine: per questo, le
    femmine durante la lattazione hanno spesso una sospensione del ciclo
    mestruale e non sono fertili;
  \item la PRL è prodotta in ugual misura anche nei maschi, ma i suoi
    effetti sono poco noti; a dosaggio fisiologico contribuisce a regolare
    la produzione di testosterone, mentre una secrezione eccessiva
    (\textbf{prolattinoma}) la blocca, causando infertilità maschile.
\end{itemize}

% ---------------------------------------------------------------------------
\section{Riepilogo degli ormoni ipofisari}

\rubrica{Adenoipofisi (lobo anteriore)}
\begin{itemize}
  \item \textbf{tireotropo} (TSH) $\rightarrow$ ghiandola tiroide:
    secrezione di ormoni tiroidei;
  \item \textbf{adrenocorticotropo} (ACTH) $\rightarrow$ corticale del
    surrene (zona fascicolata): secrezione di glucocorticoidi;
  \item \textbf{follicolo-stimolante} (FSH) $\rightarrow$ cellule
    follicolari dell'ovaio (femmina): sviluppo del follicolo e secrezione
    di estrogeni; cellule di Sertoli (maschio): spermatogenesi;
  \item \textbf{luteinizzante} (LH) $\rightarrow$ cellule follicolari
    dell'ovaio (femmina): ovulazione, formazione del corpo luteo,
    secrezione di progesterone; cellule interstiziali del testicolo
    (maschio): secrezione di testosterone;
  \item \textbf{prolattina} (PRL) $\rightarrow$ ghiandola mammaria:
    produzione di latte;
  \item \textbf{somatotropo} (GH) $\rightarrow$ tutte le cellule:
    accrescimento, sintesi proteica, mobilizzazione di lipidi e catabolismo;
  \item \textbf{melanotropo} (MSH, pars intermedia) $\rightarrow$
    melanociti: incremento della sintesi di melanina nell'epidermide.
\end{itemize}

\rubrica{Neuroipofisi (lobo posteriore)}
\begin{itemize}
  \item \textbf{vasopressina} (ADH) $\rightarrow$ tubulo contorto distale
    del rene: riassorbimento di acqua, aumento di volume e pressione del
    sangue;
  \item \textbf{ossitocina} $\rightarrow$ utero e ghiandola mammaria
    (femmina): contrazione del miometrio, eiezione di latte; dotto
    deferente e prostata (maschio): contrazione ed eiezione delle
    secrezioni.
\end{itemize}

% ---------------------------------------------------------------------------
\section{La tiroide}

La tiroide presenta \textbf{follicoli} ripieni di una sostanza ialina
chiamata \textbf{colloide}, costituita dalla proteina
\textbf{tireoglobulina}: molecole molto grandi su cui sono presenti anelli
di tirosina.

\rubrica{Sintesi degli ormoni tiroidei}
\begin{itemize}
  \item una speciale catena enzimatica lega uno o due ioni iodio alla
    tirosina, formando rispettivamente \textbf{monoiodotirosina} (MIT) e
    \textbf{diiodotirosina} (DIT);
  \item una molecola di DIT lega un MIT formando \textbf{triiodotironina}
    (T3), oppure un'altra molecola di DIT formando \textbf{tetraiodotironina}
    (tiroxina, T4);
  \item le forme secrete sono T3 e T4: T4 è la forma di gran lunga più
    abbondante, ma si trasforma in T3 nelle cellule bersaglio; T3 è
    pertanto la forma effettivamente attiva.
\end{itemize}

I \textbf{tireociti} trasportano lo iodio dai liquidi extracellulari
all'interno della cavità follicolare, dove viene coniugato con i residui
tirosinici delle molecole di tireoglobulina per formare gli ormoni
tiroidei.

\rubrica{Le cellule parafollicolari e la calcitonina}
Le \textbf{cellule parafollicolari} sono cellule cubiche che rivestono il
follicolo: pur poggiando sulla lamina basale, non raggiungono il lume del
follicolo; sono più grandi dei tireociti, poco colorabili, e sono
responsabili della produzione di \textbf{calcitonina}.

\textbf{Calcitonina}: peptide la cui azione favorisce l'accumulo di ioni
calcio nel tessuto osseo e ne aumenta l'eliminazione renale, in
antagonismo con il paratormone prodotto dalle paratiroidi.

% ---------------------------------------------------------------------------
\section{Regolazione della secrezione tiroidea}

L'ipotalamo produce un fattore, il \textbf{TRH} (ormone di rilascio
dell'ormone tireotropo), che stimola la liberazione dell'ormone ipofisario
\textbf{TSH} (ormone tiroideo stimolante). Il TSH agisce direttamente sulle
cellule follicolari tiroidee: aumenta la captazione dello iodio, stimola la
produzione di colloide e il suo riassorbimento, regolando l'immissione in
circolo di T3 e T4.

La produzione di TRH è controllata a feedback negativo dalla concentrazione
di T3 e T4: questo feedback negativo è responsabile del controllo
omeostatico del rilascio degli ormoni tiroidei.

\rubrica{Il gozzo}
La carenza di iodio nella dieta, oppure il consumo prevalente di verdure
che inibiscono il riassorbimento dello iodio (es.\ cavoli -- problema
diffuso ad esempio in alcune vallate alpine), provoca, oltre ad altri
numerosi sintomi dell'ipotiroidismo, una caratteristica ipertrofia tiroidea
che si manifesta come \textbf{gozzo ipotiroideo}: le tirosine non iodate
non possono essere secrete e si accumulano nei follicoli, continuamente
stimolate dal TSH. Esistono anche varie forme di \textbf{gozzo
ipertiroideo}, accompagnate da iperfunzione.

% ---------------------------------------------------------------------------
\section{Ipertiroidismo e ipotiroidismo}

\rubrica{Ipertiroidismo}
La sintomatologia generale comprende: ipertermia (aumento della
temperatura), dimagrimento, irritabilità e instabilità dell'umore,
insonnia, iperfagia, tachicardia con modesta ipertensione, tremore
muscolare fine.

\rubrica{Ipotiroidismo}
Nell'adulto, l'ipotiroidismo presenta sintomi opposti: ipotermia,
affaticabilità, torpore e sonnolenza, aumento di peso pur con scarso
appetito, ritenzione idrica, edema duro (\textbf{mixedema}), bradicardia e
ipotensione.

% ---------------------------------------------------------------------------
\section{Ormoni della tiroide, delle paratiroidi e del timo}

\begin{itemize}
  \item \textbf{tiroide -- tireociti}: tiroxina (T4) e triiodotironina
    (T3) $\rightarrow$ agiscono sulla maggior parte delle cellule
    corporee, con incremento dell'utilizzazione di energia, del consumo di
    ossigeno, dell'accrescimento e dello sviluppo;
  \item \textbf{tiroide -- cellule C}: calcitonina $\rightarrow$ agisce su
    ossa e reni, con diminuzione della concentrazione di ioni calcio nei
    fluidi corporei (significato incerto negli adulti in salute e non in
    gravidanza).
\end{itemize}
